\chapter{Logical-Effort Analysis}
\label{ch:logical-effort}

This chapter documents the implementation of \cref{sec:bg-le} as applied to
the top-$K$ critical paths reported by the STA engine
(\cref{sec:sta-paths}). Logical-effort analysis is purely a
\emph{post-processing} view over the STA result: it reads capacitances and
delays directly off the already-analyzed circuit and does not perform any
independent simulation of its own, so its results are always consistent
with whatever sizing assignment was current when the STA pass that produced
the critical paths was run.

\section{Per-Stage Effort Computation}
\label{sec:le-stage}

For each critical path, the sequence of gates from the driving primary
input to the endpoint is walked in order. At each stage $i$, the specific
input pin through which the path enters that gate is identified (recovered
from the same predecessor information used to reconstruct the path in
\cref{sec:sta-paths}), which determines $g_i$, the logical effort of that
particular pin, from the cell library. The on-path receiver capacitance
$C_{\mathrm{onpath},i}$ is the input pin capacitance of the next stage along
the path (or, for the final stage, the external load applied to the
endpoint's net); the off-path capacitance $C_{\mathrm{offpath},i}$ is
obtained by subtracting $C_{\mathrm{onpath},i}$ from the gate's total output
load $C_{\mathrm{load},i}$, which was already computed by the STA engine
(\cref{eq:cload}) and reflects the gate's \emph{actual, current} fan-out ---
not an idealized or path-only load. Branching effort $b_i$, electrical
effort $h_i$, and stage effort $f_i$ then follow directly from
\cref{eq:le-stage}.

\section{Path-Level Effort and the Sizing Target}
\label{sec:le-path}

Path-level quantities $G$, $B$, $H$, $F$, and $P$
(\cref{eq:le-path}) are accumulated across all stages of the path, and the
optimal per-stage effort $\hat f = F^{1/N}$ and theoretical minimum path
delay $D_{LE,\mathrm{min}}$ (\cref{eq:le-min-delay}) are computed once per
path. Every stage's actual effort $f_i$ is then compared to $\hat f$: a
stage with $f_i \gg \hat f$ is disproportionately loaded relative to its own
drive strength and is, by the theory, the best next candidate for
up-sizing. This overshoot quantity, $f_i - \hat f$, is the ranking signal
consumed by the optimizer's logical-effort-guided heuristics
(\cref{sec:opt-heuristics}).

Because each stage's desired capacitance (\cref{eq:le-target-cap}) depends
only on that stage's own current load $C_{\mathrm{total},i}$ and its own
$g_i$ --- both already known from the current, fixed sizing assignment ---
computing every stage's target capacitance does not require processing the
path in any particular order; the ``move backward along the path'' framing
used in the assignment specification is a presentational convenience,
carried over unchanged from classical hand-calculation logical-effort
sizing, rather than a genuine data dependency in this implementation.

\section{Discrete Sizing Candidates}
\label{sec:le-candidates}

The continuous target capacitance $C^{*}_{\mathrm{in},i}$
(\cref{eq:le-target-cap}) will almost never coincide exactly with one of the
three discrete sizes (\code{X1}/\code{X2}/\code{X4}) available in the
library. Rather than rounding to the single nearest size by absolute
capacitance difference, the implementation \emph{brackets} the target: it
finds the largest available cell whose input pin capacitance is still at or
below the target, and the smallest available cell whose input pin
capacitance is at or above it, and reports both as candidates (falling back
to the nearest boundary cell if the target lies outside the family's
achievable range entirely). This bracketing approach is more principled
than nearest-distance rounding because it guarantees the two candidates
straddle the theoretical ideal from both sides, which matters when the
optimizer (\cref{ch:optimization}) evaluates both bracket candidates via a
full STA re-run and lets the \emph{measured} outcome, not the continuous
approximation, determine which one is actually better once real,
non-idealized loading effects on neighboring gates are accounted for.

\begin{table}[H]
\centering
\small
\begin{tabular}{@{}lcccc@{}}
\toprule
\textbf{Quantity} & \textbf{Symbol} & \textbf{Meaning} & \textbf{Computed from} \\
\midrule
Logical effort & $g_i$ & Gate-topology cost of driving equal load & Library, per pin \\
Electrical effort & $h_i$ & On-path fan-out ratio & \cref{eq:le-stage} \\
Branching effort & $b_i$ & Off-path loading tax & \cref{eq:le-stage} \\
Stage effort & $f_i$ & $g_i b_i h_i$ & \cref{eq:le-stage} \\
Path effort & $F$ & $\prod_i f_i$ & \cref{eq:le-path} \\
Optimal stage effort & $\hat f$ & $F^{1/N}$ & \cref{eq:le-min-delay} \\
Theoretical min.\ delay & $D_{LE,\mathrm{min}}$ & $\tau(NF^{1/N}+P)$ & \cref{eq:le-min-delay} \\
Target capacitance & $C^{*}_{\mathrm{in},i}$ & Ideal input cap.\ for stage $i$ & \cref{eq:le-target-cap} \\
\bottomrule
\end{tabular}
\caption{Summary of quantities reported per stage/path in
\texttt{logical\_effort\_paths.csv}, \texttt{logical\_effort\_stages.csv},
and \texttt{logical\_effort\_candidates.csv}.}
\label{tab:le-summary}
\end{table}
